\section{Objetivos de la Investigación}
\label{chap:objetivos}

\subsection{Objetivo General}
\label{sec:objetivo_general}
Construir y evaluar un \textbf{bridge universal extensible de simulación para APIs nativas de Capacitor} en aplicaciones Ionic, con \textbf{procedimiento de cobertura en Node.js} (Jest y/o Jasmine) y \textbf{capa auxiliar MCP}, contrastando resultados con una \textbf{línea base} documentada en el mismo proyecto de referencia.

Responde al problema general de la sección~\ref{sec:problema_general}.

\subsection{Objetivos Específicos}
\label{sec:objetivos_especificos}
Los objetivos O1--O7 responden a los problemas P1--P7 (sección~\ref{sec:problemas_especificos}). La numeración es paralela; la redacción de cada ítem varía según el entregable.

\begin{enumerate}[label=\textbf{O\arabic*.}, leftmargin=*]
    \item \label{obj:simulador}
    \textbf{[P\ref{prob:simulador}]} Desarrollar en TypeScript el bridge con perfiles configurables para \tcode{geolocation}, \tcode{camera} y \tcode{preferences}, cubriendo éxito, error, permisos denegados y valores límite según el contrato publicado de cada plugin.

    \item \label{obj:integracion}
    Integrar el bridge en al menos un proyecto Ionic de referencia: preset Jest, \tcode{setupFilesAfterEnv} y mapeo de módulos documentados; en Jasmine/Angular, el equivalente razonable. Activar \textit{fallback} con \tcode{Proxy} para \tcode{@capacitor/*} fuera del núcleo. \textbf{[P\ref{prob:integracion}]}

    \item \label{obj:cobertura}
    \textbf{[P\ref{prob:cobertura}]} Aplicar un método de cobertura: inventariar usos de los tres plugins, listar escenarios mínimos por rama y escribir o ajustar tests hasta elevar líneas y ramas en \tcode{jest --coverage}, sin binarios en teléfono.

    \item \label{obj:mcp}
    Implementar servidor MCP auxiliar con herramientas acotadas (lanzar suite con cobertura, leer agregados del reporte, Sonar solo lectura si existe). Dejar por escrito permisos, rutas y qué no hará el servidor. \textbf{[P\ref{prob:mcp}]}

    \item \label{obj:mcp_valor}
    \textbf{[P\ref{prob:mcp_valor}]} Comparar, en el mismo subconjunto de corridas, flujo MCP asistido frente a script o CI lineal: medir $T_{diag}$, pasos manuales y si el siguiente perfil de simulación se elige con la misma criterio.

    \item \label{obj:metricas}
    Fijar métricas y diseño $B \rightarrow P$: cobertura $C$, tiempo de preparación de mocks $T$, pasaje de suite; tabla de caracterización de la línea base (8+ componentes, 2\,500+ LOC, 6+ mocks previos). \textbf{[P\ref{prob:metricas}]}

    \item \label{obj:transferencia}
    \textbf{[P\ref{prob:transferencia}]} Redactar guías cortas para extender perfiles a otros plugins y enlazar el flujo en CI, con límites explícitos del alcance actual.
\end{enumerate}
